\section{An exact nonsplit Gaussian $E_8$ twelve-branch source}
\label{sec:im-e8-tuple}

\paragraph{Status and scope.}
The assertions in D-E8T-01--03 below are established by the displayed
arguments and the exact certificate
\texttt{agents/imprimitive/certify\_e8\_tuple.py}, with full output in
\texttt{agents/imprimitive/e8\_search/tuple\_certificate.json}.
They establish product-one generation, the actual Weil signature and
polarization class, and a boundary with toric rank two. They do not
assert rank nine of the period map or algebraicity of its Weil classes.
The tuple was found independently by an exhaustive conjugacy-orbit
search; no tuple, cup matrix, or rank theorem from the unrefereed
Gaussian $E_8$ manuscript is a premise.

\subsection*{D-E8T-01. An explicit generating tuple}
Use the simple-root basis and Cartan matrix $H$ of HD-E8-01. Thus
$s_a=I-e_aH_{a,*}$, and the edges of the Dynkin diagram are
$(1,3),(3,4),(2,4),(4,5),(5,6),(6,7),(7,8)$.
Let $j_0$ be the following integral matrix:
\[
j_0=\begin{pmatrix}
0&1&1&0&-1&0&-1&0\\
1&2&0&0&-1&0&-1&-1\\
0&3&1&0&-2&0&-1&-1\\
1&4&1&0&-2&-1&-1&-2\\
1&3&1&0&-2&0&-1&-2\\
1&2&0&1&-2&0&-1&-1\\
1&2&0&0&-1&0&-1&0\\
0&1&0&0&0&0&-1&0
\end{pmatrix}.
\]
This is the representative of $2D_4(a_1)$ reconstructed as an explicit
simple-reflection word in HD-E8-01; its word, matrix, and provenance
are retained in \texttt{research/e8\_class\_exploration.json}.
Put
\[
u=s_6s_7s_8s_1,\qquad j=uj_0u^{-1},\qquad
c=s_1s_2s_3s_4s_5s_6,\qquad w=jc^{-1},\qquad h=ij,\quad g=iw.
\]
The twelve-entry tuple is
\[
\boldsymbol t=(s_1,s_1,s_2,s_2,s_1,s_2,s_3,s_4,s_5,s_6,h,g).
\tag{E8T}
\]
It has product one, generates $G=\mu_4W(E_8)$, and has branch orders
$(2^{11},12)$.

\begin{proof}
The matrices satisfy $j^2=-I$ and $j^tHj=H$. Consequently
$h^2=I$ and
\[
hg=-jw=-j^2c^{-1}=c^{-1}.
\]
The two initial pairs cancel, so the product of (E8T) is $chg=I$.
The exact characteristic polynomial of $w$ is
\[
(x-1)^2\Phi_3(x)\Phi_{12}(x)
=x^8-x^7-x^6+2x^4-x^2-x+1.
\]
The eight eigenangle fractions of $g=iw$, in increasing order, are
\[
\frac16,\ \frac14,\ \frac14,\ \frac13,\ \frac7{12},\
\frac23,\ \frac56,\ \frac{11}{12}.
\tag{E8T-spectrum}
\]
They are all nonzero, sum to four, and give order twelve.
The certificate independently checks $g^k\ne I$ for $1\le k<12$
and $g^{12}=I$.

Here is a certificate of generation which does not enumerate the
$1393459200$ elements of $G$. Starting with the six roots
$e_1,\ldots,e_6$, close under the seven transformations
$s_1,\ldots,s_6,j$. The resulting set has 240 distinct vectors, each
of $H$-norm two, and equals the set obtained by closing all eight
simple roots under all eight simple reflections. For every new root
the output records its parent root and the transformation applied;
these are finite exact integer equalities. The latter set is the full
$E_8$ root system. In the group generated by the tail
$(s_1,\ldots,s_6,h,g)$, conjugation by $h=ij$ carries the reflection
in a root $v$ to the reflection in $jv$: the scalar $i$ cancels in
conjugation. Hence the root-closure certificate implies that this
group contains every $E_8$ root reflection and therefore contains
$W(E_8)$. It contains $j$, and hence $hj^{-1}=iI$ as well. Thus even
the eight-entry tail generates all of $G$.

For reproducibility, the search first closed $j_0$ under conjugation
by all eight simple reflections and found precisely 15120 distinct
matrices, verifying $j^2=-I$ for each. For $c=s_1\cdots s_6$, exactly
1080 of these gave the characteristic polynomial of class 57 or 58
for $jc^{-1}$. The selected matrix is orbit index 179 and is obtained
by successive conjugations $[1,8,7,6]$. It is the first such match
examined for which the stated root closure gives 240 roots.
\end{proof}

\subsection*{D-E8T-02. The actual sixfold and its pairing}
A connected $G$-cover of $\mathbf P^1$ with tuple (E8T) has genus
$3077222401$. Its natural-representation primitive factor is an
abelian sixfold of Weil type for $K=\mathbf Q(i)$, of signature
$(3,3)$, with raw Hermitian determinant class $[-3]$. Its Hurwitz
component has dimension nine.

\begin{proof}
The $E_8$ reflection representation is absolutely irreducible: the
simple roots span, and their connected nonorthogonality graph forces
any orthogonal decomposition into invariant complex subspaces to
place all roots in one summand. The matrices are defined over $K$,
and the scalar $iI$ has trace $8i$, so the character field is exactly
$K$. The associated rational simple block is $\operatorname{Mat}_8(K)$.
A rational Rosati-fixed primitive projector in this block therefore
cuts out one copy of the multiplicity factor as an actual polarized
abelian subvariety of the covering Jacobian. Equivalently, project
orthogonally to a root line inside the central natural block. The
central scalar $iI$ preserves its integral homology lattice.

There are ten rank-one reflection entries of age $1/2$, one involution
$h$ with four eigenvalues $-1$ and four eigenvalues $1$, hence age two,
and the full-moving-rank entry $g$ with age four by
(E8T-spectrum). Their total moving rank is $10+4+8=22$ and their
total age is eleven. Chevalley--Weil gives
\[
\dim_K H^1_{\mathrm{par}}=-16+22=6,\qquad
h^{1,0}_{\sigma}=h^{1,0}_{\bar\sigma}=11-8=3.
\]
Thus the primitive factor has dimension six and the asserted
signature. Riemann existence and the usual ordered Hurwitz space
supply a finite etale covering of an open part of $M_{0,12}$, so its
component has dimension nine. Riemann--Hurwitz gives
\[
g(C)=1+\frac{1393459200}{2}
\left(-2+11\frac12+\frac{11}{12}\right)=3077222401.
\]

For the actual polarization, use the parabolic cup construction of
D-DET-01. To make the verification independently reproducible, for
each inertia matrix $t_a$ take the first independent columns $U_a$
of $I-t_a$, let $T_a$ be its restriction to that moving space, and put
$Q_a=t_1\cdots t_{a-1}$. The relation matrix has blocks $Q_aU_a$.
The coboundary matrix has blocks consisting of the coordinates of
$(I-t_a)x$ in $U_a$. The upper triangular Seifert matrix has blocks
\[
S_{ab}=(Q_aU_a)^*H(Q_bU_b)\quad(a<b),\qquad
S_{aa}=(I-T_a)^{-*}U_a^*HU_a.
\]
It satisfies, exactly,
\[
S+S^*=R^*HR,\qquad B^*S=HR,\qquad RB=0,
\qquad \operatorname{rank}R=\operatorname{rank}B=8.
\]
Thus its restriction to $\ker R$ descends to the anti-Hermitian cup
form on the quotient, up to the common cup-convention sign, which
does not affect its determinant in even dimension. The output saves
$R,B,S$ and every column of the quotient basis. The resulting matrix is
\[
S_{\mathrm{par}}=\begin{pmatrix}
0&0&-4&0&2-4i&2\\
0&0&0&-4&2-2i&0\\
4&0&0&0&-2+4i&-2\\
0&4&0&0&-2+2i&0\\
-2-4i&-2-2i&2+4i&2+2i&-6i&3-3i\\
-2&0&2&0&-3-3i&-2i
\end{pmatrix},
\qquad \det S_{\mathrm{par}}=1536.
\]
Converting to a Hermitian form by multiplying by $i$ changes the
six-dimensional determinant to $-1536$. Since $1536=3\cdot2^9$
and $2$ is a Gaussian norm, its raw determinant class is $[-3]$.
This is not the split class $[-1]$, since the valuation at the
inert prime three is odd. This cup form is the multiplicity
polarization; changing the chosen primitive root-line factor scales
it by a positive rational number, whose sixth power is a norm.
\end{proof}

\subsection*{D-E8T-03. A rank-two toric boundary and its anisotropic residual factor}
Degenerate the two initial equal-reflection pairs in (E8T) onto
separate rational bubbles. The natural multiplicity space of the
dual graph has $K$-dimension two. The remaining polarized
weight-one factor has signature $(1,1)$, with anti-Hermitian matrix
\[
S_{\mathrm{res}}=
\begin{pmatrix}-6i&3-3i\\-3-3i&-2i\end{pmatrix},
\qquad \det S_{\mathrm{res}}=6.
\]
It is anisotropic over $K$. The positive sum of the two geometric
node logarithms has rank two on the natural multiplicity space.
In particular the actual six-dimensional $K$-Hermitian space has
Witt index exactly two.

\begin{proof}
Each bubble has inertia tuple $(s_a,s_a,1)$; its connected covering
components are rational double covers of $\mathbf P^1$. The central
component is connected, because its tail tuple
$(s_1,\ldots,s_6,h,g)$ generates $G$ by D-E8T-01. The attachment
inertia is trivial. Hence the two sets of edges have permutation
modules $K[G]$, the bubble vertices have modules
$K[G/\langle s_1\rangle]$ and $K[G/\langle s_2\rangle]$, and the
central vertex contributes the trivial module. Exactness of the
connected graph's incidence sequence gives, in its representation
ring,
\[
H_1(\Gamma,K)=2K[G]-K[G/\langle s_1\rangle]
                         -K[G/\langle s_2\rangle].
\]
The natural representation has multiplicity eight in $K[G]$ and
multiplicity seven in each induced trivial module. Its graph
multiplicity is therefore $16-7-7=2$.

The central curve has the eight-entry tail tuple. Its total moving
rank is $6+4+8=18$, so its parabolic multiplicity dimension is
$18-16=2$; its two holomorphic multiplicities are $3+2+4-8=1$.
Its genus is
\[
1+\frac{1393459200}{2}
 \left(-2+7\frac12+\frac{11}{12}\right)=1683763201.
\]
The same exact cup computation as above gives the displayed residual
matrix. The Hermitian form $iS_{\mathrm{res}}$ has first diagonal
entry six and determinant $-6$; completing the square diagonalizes
it over $K$ to $\langle6,-1\rangle$. An isotropic vector would make
six a norm from $K$, impossible because its three-adic valuation is
odd. Thus the residual plane is anisotropic.

For an equivariant semistable smoothing, the vanishing cycles of
the nodes have precisely the vertex-incidence relations; their
natural multiplicity has dimension two by the graph calculation.
A positive smoothing logarithm is a sum of positive multiples of
the Picard--Lefschetz rank-one tensors of these cycles. The positive
coefficients ensure that its rank is the dimension of their span,
namely two over $K$. Its image is isotropic for the actual
polarization. The weight-one graded quotient is the central
curve's residual factor just computed. It follows that the Witt
decomposition has two hyperbolic planes and this anisotropic plane.

The two node parameters alone do not establish rank nine. In the
unitary domain, a boundary with isotropic dimension two has four
toric-block tangent directions, four mixed extension directions,
and one residual-period direction. A proof of dominance must
control the actual image in all nine directions, or give a separate
full-rank codifferential calculation at a smooth cover. This remains
separate from the exact tuple and pairing results proved here.
\end{proof}
